\chapter{KULLANILMASI PLANLANAN YÖNTEMLER VE TEKNOLOJİLER}
\thispagestyle{empty}

\section{Yöntemler}

\textbf{Veri Seti İnşası -- Custom-WikiCS:}
Projede kullanılan veri seti, mevcut Wiki-CS \cite{mernyei2020wikics} kıyaslama veri setinin topolojisi temel alınarak yeniden inşa edilmiştir. Süreç şu adımlardan oluşmaktadır:
\begin{enumerate}
    \item Wiki-CS veri setinin graf topolojisi (11.701 düğüm, 291.039 yönlü kenar) alınmıştır.
    \item Düğüm başlıkları işlenerek Wikipedia'daki gerçek makale başlıklarıyla eşleştirilmiştir.
    \item Kaggle üzerinde bulunan ``Wikipedia STEM Corpus'' veri seti kullanılarak eşleşen makalelerin metinleri çıkarılmıştır.
    \item Corpus'ta bulunmayan makalelerin metinleri Wikipedia API aracılığıyla tamamlanmıştır.
    \item Elde edilen metinler, iki farklı cümle gömme modeli ile vektörleştirilerek düğüm öznitelikleri oluşturulmuştur.
\end{enumerate}

Bu yaklaşım, orijinal Wiki-CS'nin yalnızca GloVe ortalaması kullanan basit metin temsilini modern anlamsal gömme vektörleriyle değiştirirken, graf yapısının güvenilirliğini korumaktadır.

\textbf{Cümle Gömme Modelleri:}
Makale metinlerinin vektör temsilini oluşturmak için iki farklı ön eğitimli model karşılaştırmalı olarak değerlendirilmektedir:
\begin{itemize}
    \item \textbf{BAAI/bge-m3:} 1024 boyutlu gömme vektörleri üreten, çok dilli ve çok görevli bir model.
    \item \textbf{Alibaba-NLP/gte-modernbert-base:} 768 boyutlu gömme vektörleri üreten, ModernBERT mimarisine dayalı bir model.
\end{itemize}
Hangi modelin GNN ile birleştirildiğinde daha iyi öneri performansı sağladığını belirlemek amacıyla deneyler sürdürülmektedir.

\textbf{Graf Sinir Ağı -- GAT:}
Makaleler arası yapısal ilişkileri öğrenmek için Graph Attention Networks (GAT) \cite{velickovic2018graph} kullanılacaktır. GAT, her komşu düğüme farklı dikkat ağırlıkları atayarak hangi bağlantının daha önemli olduğunu öğrenebilmektedir. Temel karşılaştırma modeli (baseline) olarak GCN \cite{kipf2017semi} kullanılacaktır.

\textbf{Büyük Dil Modeli -- Hiyerarşik Sunum:}
GNN tabanlı öneri sonuçlarının kullanıcıya hiyerarşik bir öğrenme patikası olarak sunulması için yerel olarak çalıştırılabilen açık kaynaklı bir Büyük Dil Modeli (örn. DeepSeek-R1 veya benzer bir model) kullanılması planlanmaktadır. Model seçiminde 8~GB VRAM sınırına uyumluluk gözetilmektedir.

\section{Yazılım}

\begin{itemize}
    \item \textbf{Python:} Temel programlama dili.
    \item \textbf{PyTorch \& PyTorch Geometric (PyG):} GNN modellerinin (GAT, GCN) oluşturulması ve eğitimi için.
    \item \textbf{HuggingFace Transformers / Sentence-Transformers:} Cümle gömme modellerinin çalıştırılması için.
    \item \textbf{NetworkX / iGraph:} Graf analizi ve görselleştirme için.
    \item \textbf{Ollama:} Yerel LLM kullanımı için.
\end{itemize}

\section{Donanım}

Deneyler, GPU hızlandırmalı bir dizüstü bilgisayar üzerinde yürütülmektedir:
\begin{itemize}
    \item \textbf{GPU:} NVIDIA RTX 4070 Laptop (8~GB VRAM) -- model eğitimi ve çıkarım için.
    \item \textbf{CPU:} AMD Ryzen 9 8945HS (4~GHz) -- veri ön işleme ve graf analizi için.
    \item \textbf{RAM:} 32~GB -- büyük graf yapılarının bellekte tutulması için.
\end{itemize}